\section{Data and Methodology}

\subsection{Data and Sample Construction}

Our empirical analysis rests on a monthly panel of index-level Price-to-Book (P/B) ratios spanning January 2004 to April 2026. The core universe comprises four benchmark indices: the Korea Composite Stock Price Index (KOSPI), the Tokyo Stock Price Index (TOPIX), the S\&P 500, and the MSCI Emerging Markets index. These four series constitute the primary study universe and support both descriptive and event-study analysis through December 2024, which we treat as the principal study endpoint consistent with the paper's pre-registered scope. A follow-on analysis window extending through April 2026 supports the Korea Value-Up event-study and robustness specifications, which require additional post-event observation months given the recency of Korea's 2024 reform milestones.

The regional robustness analysis draws on a broader set of MSCI country and regional indices---MSCI Hong Kong, MSCI Taiwan, MSCI EM Asia, MSCI China, MSCI India, and MSCI Indonesia---which serve two distinct purposes. In the Korea alternative-comparator robustness specification, MSCI EM Asia replaces TOPIX as the denominator of the Korea valuation spread, allowing us to assess whether Japan-specific post-reform appreciation mechanically contaminates our headline Korea--Japan spread estimate. In the synthetic control analysis, the same regional indices constitute the donor pool against which a counterfactual Japan is constructed.

All P/B series are measured at the aggregate index level and are available at monthly frequency. Index-level P/B ratios have several methodological advantages over the firm-level multiples that dominate the microeconometric governance literature reviewed above. Because they reflect the value-weighted composition of each market's listed equity, they integrate information across thousands of firms and smooth idiosyncratic cross-sectional noise. This property is particularly important for identifying structural valuation wedges---such as the Korea Discount---that persist across sectors, size groups, and firm-level governance variation. The tradeoff is that index-level analysis precludes controlling for firm-level observables such as profitability, leverage, or R\&D intensity; our identification strategy therefore relies on the plausible assumption that such compositional differences between KOSPI and TOPIX are approximately stable within event windows of the length we employ (12 months pre-event to 24 months post-event).

\subsection{Descriptive Analysis}

Figure 1 presents the 20-year P/B time series for KOSPI, TOPIX, S\&P 500, and MSCI EM over the 2004--2024 study window, with Japan's three primary governance reform dates annotated by vertical reference lines. This figure establishes the key stylized facts motivating our event-study design. The KOSPI--TOPIX gap, visible at the start of the sample and persistent throughout, narrows modestly during the global financial crisis and its aftermath, then resumes a downward trajectory by the mid-2010s. The most striking feature of the descriptive record is the divergent trajectory of the two indices following Japan's March 2023 TSE P/B reform mandate, a divergence that the event-study analysis formalizes.

\subsection{Primary Spread and Event-Study Design}

The primary outcome variable throughout our analysis is the KOSPI--TOPIX P/B spread: the monthly difference between the KOSPI index-level P/B ratio and the TOPIX index-level P/B ratio. Choosing Japan as the comparator is motivated by structural, geographic, and institutional similarities that make it the most credible counterfactual available. Japan and Korea are both capital-account-open, export-oriented economies with concentrated corporate ownership structures, and both listed equity markets are classified as developed by major index providers. Crucially, Japan has passed through a staged sequence of governance reforms over the past decade, culminating in a hard exchange-level mandate in 2023---precisely the reform trajectory that Korea's policymakers and the literature argue Korea must replicate. These features justify Japan as the benchmark against which Korean governance policy can be evaluated.

We use a stacked event-study design to estimate the cumulative abnormal spread around each reform event. The stacking approach, which constructs a separate cohort dataset for each event date and then pools them into a single regression, is well-suited to settings with multiple non-overlapping or partially overlapping reform events, because it allows each cohort's pre-event trend to serve as its own counterfactual baseline without constraining effects to be homogeneous across events \citep{cengiz_stacked_did_2019}.

For each event date $t_k$ ($k = 1, \ldots, K$), we construct a cohort panel by extracting all months within a symmetric window around that event. The stacking window for the Japan analysis runs from 36 months before to 24 months after each event date, providing a pre-event observation span sufficient to estimate a reliable baseline trend. Months within each cohort are indexed by their calendar-month distance from $t_k$, denoted $\tau \in [-36, +24]$, where negative values are pre-event and positive values are post-event. For the Korea analysis, the stacking window is $[-36, +20]$ months, with the post window truncated to 20 months to respect the available post-reform observation period through April 2026.

\paragraph{Abnormal Spread Construction.} Within each cohort, we estimate a pre-event linear trend in the KOSPI--TOPIX spread over the pre-period $\tau \in [-36, -1]$, fitting a simple OLS regression of the spread on $\tau$. This linear trend projection provides the expected spread for all months in $[-36, +24]$, including the post-event months. The abnormal spread for each month is then defined as the realized spread minus the expected spread. Normalizing the baseline period at $\tau = -1$ (the month immediately before the reform event) ensures that estimated effects capture deviations from the pre-event trajectory rather than level differences in the pre-period mean, and removes transient level shifts that may have accumulated at event announcement.

\paragraph{Stacked Regression and Cumulative Abnormal Returns.} We pool the cohort-level abnormal spread panels into a stacked dataset and estimate a saturated cohort-by-relative-time OLS model:
\begin{equation}
    \text{AbnSpread}_{k,\tau} = \sum_{k=1}^{K} \sum_{\tau \neq -1} \beta_{k,\tau} \cdot \mathbf{1}[\text{cohort} = k, \text{event\_rel\_time} = \tau] + \varepsilon_{k,\tau}
\end{equation}
where the base category is $\tau = -1$ for each cohort, so $\beta_{k,-1} \equiv 0$ by construction. The saturated design is fully flexible: it imposes no parametric restriction on the pattern of abnormal spreads within or across cohorts, and makes no restrictions on treatment effect heterogeneity across events. The cumulative abnormal return (CAR) for cohort $k$ through month $\tau^*$ is:
\begin{equation}
    \text{CAR}_{k,\tau^*} = \sum_{\tau=-12}^{\tau^*} \hat{\beta}_{k,\tau}
\end{equation}
where the event window for the reported CARs begins at $\tau = -12$ rather than the full stacking window start at $\tau = -36$, ensuring that CARs reflect the policy-relevant window around each reform date rather than the broader pre-event period used for baseline estimation.

Because the saturated cohort-by-relative-time design exhausts all degrees of freedom in the time dimension within each cohort, the design matrix is not full rank when observations from multiple cohorts overlap in calendar time---as they do in the Korea analysis, where three Value-Up reform dates cluster within a seven-month window in 2024. In this setting, robust sandwich standard errors are undefined for the saturated design, and we therefore report estimated coefficients and CARs as descriptive statistics rather than as estimates with formal inferential claims. This is the appropriate treatment in an exploratory context where overlapping event windows are a design feature to be disclosed rather than a source of bias to be eliminated.

\subsection{Japan Event-Study: Reform Benchmark Identification}

The Japan event-study uses three governance reform events: (i) the Financial Services Agency Stewardship Code (February 2014), which established behavioral expectations for institutional investors; (ii) the Tokyo Stock Exchange Corporate Governance Code (June 2015), which imposed formal board composition and disclosure standards on listed companies; and (iii) the TSE P/B Reform Request (March 2023), under which the TSE required all listed firms with P/B below 1.0 to disclose capital efficiency improvement plans or explain non-compliance. These three events span Japan's full governance reform arc from soft institutional guidance to hard exchange-level mandate, and they constitute the governance sequence that Korea's policy trajectory is being benchmarked against.

The direction of the spread effect depends on whether the reform in question caused TOPIX to appreciate or KOSPI to depreciate (or both) in relative terms. A persistent negative CAR in the KOSPI--TOPIX spread following a Japanese reform would indicate that Japan's valuation improved relative to Korea's---consistent with the governance channel if the reform was effective in raising TOPIX values but did not produce a symmetric response in KOSPI. A null or positive CAR would indicate no widening of the Korea Discount in response to Japan's reform milestones.

\subsection{Geopolitical Risk Falsification}

The literature review established that North Korean geopolitical risk is a commonly cited but empirically weak driver of the Korea Discount. We test this formally using the Geopolitical Risk index for Korea (GPRC\_KOR) constructed by \citet{caldara_gprc_2022}, which provides a monthly country-level measure of media-reported geopolitical threats. To operationalize the escalation hypothesis in its strongest form, we construct a binary escalation dummy equal to one when GPRC\_KOR exceeds its 75th percentile over the 2004--2024 study window---a threshold that identifies the most severe geopolitical risk episodes rather than the chronic background level.

We estimate the following cross-sectional time-series model:
\begin{equation}
    \text{KOSPI\_PB}_t = \alpha + \beta_1 \cdot \text{GPR\_Escalation}_t + \beta_2 \cdot \text{TOPIX\_PB}_t + \sum_{y} \gamma_y \cdot \mathbf{1}[\text{year} = y] + \varepsilon_t
\end{equation}
where the TOPIX P/B control absorbs common global equity market movements, and year fixed effects absorb secular trends in both series over the two-decade panel. Standard errors are HC3 heteroskedasticity-robust. The key parameter is $\beta_1$: if geopolitical risk is a meaningful driver of the Korea Discount, KOSPI P/B should be systematically lower in high-escalation months even after controlling for TOPIX co-movement and year-level trends. A precisely estimated zero would undermine the geopolitical risk hypothesis and support the governance-centric interpretation pursued throughout this paper.

We note that $\hat{\beta}_1$ is a partial identification result---the GPR escalation dummy may absorb global risk-off episodes that simultaneously depress multiple equity markets, in which case the coefficient is an upward-biased estimate of the pure North Korea threat effect. Our preferred interpretation of a near-zero coefficient is therefore that even this potentially attenuated estimate of the geopolitical risk channel is negligible.

\subsection{Synthetic Control Validation}

As an alternative identification strategy for the 2023 TSE P/B reform, we apply the synthetic control method of \citet{abadie_synthetic_control_2010} to construct a counterfactual Japan P/B trajectory. The treatment unit is the TOPIX (Japan) index, and the donor pool consists of seven MSCI regional and country indices that did not undergo an equivalent P/B governance mandate in 2023: MSCI Hong Kong, MSCI Taiwan, MSCI Emerging Markets, MSCI EM Asia, MSCI China, MSCI India, and MSCI Indonesia. South Korea (KOSPI) is excluded from the donor pool to preserve the independence of the primary KOSPI--TOPIX spread analysis. All donor countries are regional or emerging market peers with no exchange-level P/B mandate reform during the study period, satisfying the SUTVA assumption that donor units are unaffected by the Japanese treatment.

Japan's absolute P/B level lies structurally below several donor markets, creating a convex hull problem under which standard synthetic control matching on levels would concentrate all weight on the single lowest-valued donor. To resolve this, we demean each series to its pre-treatment sample mean before optimization, allowing the synthetic control weights to match on trend shape rather than levels. The pre-treatment period for weight estimation spans January 2004 to February 2023 (19 years), providing a long baseline over which global governance and ESG trends common to all markets are averaged out.

The synthetic Japan gap series---realized TOPIX P/B minus weighted synthetic TOPIX P/B---is our counterfactual outcome for the post-2023 period. We assess the reliability of the synthetic control estimate through two placebo tests. The in-time placebo assigns a fake treatment date of January 2019 and re-runs the full estimation on pre-actual-reform data, verifying that the pre-treatment fit holds and that no spurious post-``treatment'' gap emerges before the true reform. The in-space placebo iterates over each donor market, treating each as the treated unit with the remaining donors plus Japan constituting the placebo pool, and reports the distribution of placebo gap series. A Japan gap that is visually distinguishable from the placebo distribution in the post-2023 period constitutes informal evidence that the observed gap is attributable to the TSE reform rather than to structural differences between Japan and its regional peers.

\subsection{Korea Value-Up Event Study}

We apply the same stacked event-study design to three Korean governance reform milestones from Korea's 2024 ``Corporate Value-Up'' program: (i) the program launch (February 26, 2024), which announced the policy framework; (ii) the guidelines unveiling (May 2, 2024), which released implementation guidelines for voluntary firm participation; and (iii) the implementation push (August 12, 2024), which reiterated regulatory pressure for firms to file value enhancement plans. These dates are locked in \texttt{config.py} and imported by all analysis scripts to prevent any post-hoc selection of event dates.

The Korea event-study uses the same stacking window parameters as the Japan analysis ($[-36, +20]$ months), with the post window capped at 20 months given data availability through April 2026. The primary spread is KOSPI minus TOPIX P/B, preserving strict comparability with the Japan event-study. An alternative spread using KOSPI minus MSCI EM Asia P/B is reported as a robustness check: because MSCI EM Asia includes Korea as a constituent, this benchmark captures Korea's valuation relative to the broader Asian emerging market complex without depending on the Japan-specific post-2023 reform dynamics that drive the primary spread.

The three Korea reform dates cluster within a seven-month window (February, May, and August 2024), creating calendar-time overlap across cohort event windows. Under overlap, months in early 2024 simultaneously contribute to the post-event segment of the February cohort and the pre-event segment of the May cohort. We retain and explicitly annotate these overlapping observations rather than dropping them or restricting the event set, for two reasons. First, the Korea Value-Up policy milestones are analytically distinct, and collapsing them into a single event would discard informative variation in market responses across program stages. Second, since we report descriptive CARs rather than inferential estimates, the overlap does not distort the coefficient interpretation in the same way it would distort a difference-in-differences variance estimator.

\subsection{Counterfactual Projection}

Figure 4 extends the synthetic control analysis to a forward-looking scenario. Using the average monthly lift in the TOPIX--synthetic TOPIX gap during months 1--18 post-reform as an estimate of the monthly valuation increment attributable to Japan's governance mandate, we apply this lift to KOSPI's last observed P/B level (December 2024) to project Korea's hypothetical P/B trajectory over a 60-month horizon under the assumption that Korea implements an analogous exchange-level reform. Uncertainty is quantified by the pre-treatment root mean square prediction error (RMSPE) from the synthetic control estimation, which reflects the quality of fit in the pre-treatment period and provides a conservative model uncertainty band.

This projection is explicitly illustrative rather than a causal forecast. The lift estimate is drawn from Japan's experience under a specific institutional context, and Korea's governance trajectory, Chaebol composition, and investor base differ in ways that would affect the magnitude and timing of any valuation response. The projection's purpose is to calibrate the order of magnitude of potential gains under structural reform, not to provide a precise point estimate.

\subsection{Robustness Specifications}

We report four robustness checks for the Korea event-study. First, the \textit{narrow 2024 rollout} specification is the primary baseline, using the three February--August 2024 dates with a 20-month post window. Second, a \textit{window sensitivity} specification truncates the post window to 12 months, testing whether the headline result is sensitive to how much of the follow-on period is included given the limited post-event history. Third, the \textit{spaced follow-through} specification replaces the clustered 2024 dates with a temporally spread event sequence---the February 2024 program launch, the July 2025 mandatory governance disclosure expansion, and the February 2026 Value-Up dividend disclosure rule---allowing each event's effects to propagate without window overlap, though at the cost of a maximum post window of only 2 months for the latest event given the current study endpoint. Fourth, the Korea--Asia alternative comparator specification, discussed above, replaces TOPIX with MSCI EM Asia as the spread denominator.

Together, the primary result and these robustness specifications map out the sensitivity of our conclusions to comparator choice, event date selection, and post-event window length---the three most consequential design choices in an index-level event study of this kind.
